\section{Introduction}
\label{sec:intro}

\subsection{Motivation}
\label{sec:intro:motivation}

Between April 2025 and the end of 2025-Q1 2026, the inter-agent
interoperability landscape passed through what is, by any reasonable
counting, a Cambrian phase. Google launched the Agent2Agent (A2A)
protocol in April 2025 and donated it to the Linux Foundation later in
the year~\cite{spec_a2a, lf_aaif}. Anthropic's Model Context Protocol
(MCP), introduced in late 2024, became the de-facto tool-binding
standard and was adopted by every major model vendor within
months~\cite{spec_mcp}. The Agent Network Protocol (ANP) published its
white paper in August 2025~\cite{spec_anp}. IBM's ACP/BeeAI,
Cisco-led AGNTCY and its SLIM secure messaging layer, Coral Protocol,
the LOKA framework, and the Agent Capability Negotiation and Binding
Protocol (ACNBP) all reached public specifications within the same
twelve-month window~\cite{spec_coral, spec_loka, spec_acnbp}.
In parallel, the agentic-commerce stack took shape: Google's Agent
Payments Protocol (AP2), Visa's Trusted Agent Protocol (TAP),
Mastercard's Agent Pay, and the Ethereum community's ERC-8004.
A fifth specification, the Universal Commerce Protocol (UCP),
followed at the start of 2026 and covers the commerce transaction
around the payment rather than the payment itself.
These specifications introduced a third architectural layer. It is
concerned not with tool invocation or peer messaging but with
cryptographically-attested commerce settlement on behalf of an
agent principal.

Four broad surveys of inter-agent protocols already
exist~\cite{yang2025survey, ehtesham2025interop, kong2025communication,
commcentric2025survey}. They establish taxonomies and identify research
directions. Kong et al.~\cite{kong2025communication} additionally
provide a deep security analysis in 48 pages and 18 figures. None,
however, proposes a normalized scoring rubric that could be applied
uniformly across the protocol landscape; none covers the
agentic-commerce stack at the spec level; none maps industry adoption
by vertical; and none ships a machine-readable companion artifact that
remains usable as protocol versions drift during the typical 6 to 9
month review cycle of an archival journal. The contribution gap
motivating this survey is
precisely that intersection:
\emph{a rubric that scores existing specs uniformly, applied across the
full inter-agent protocol landscape including commerce, grounded in
named industry deployments, and re-runnable from a machine-readable
score file.} Re-runnable is a claim about the artifact, not about
scoring reliability: every score here is one rater's, and
\S\ref{sec:rubric:calibration} states what agreement evidence that
does and does not give us.

The closest precedent in method predates these surveys. Chopra,
Christie, and Singh define criteria for evaluating communication
protocol languages and compare five academic languages against
them~\cite{chopra2020protocol}. Those languages are formal notations
for specifying multiagent interactions. The protocols surveyed here
are industry specifications that deployed LLM agents use. They raise
questions of identity, authorization, and governance that lay outside
that study's scope, and the rubric of \S\ref{sec:rubric} adds axes
for them.

\subsection{Why now: the 2024 to 2026 protocol explosion}
\label{sec:intro:whynow}

Three forces converged to produce the protocol explosion that this
survey records. First, the maturation of frontier LLMs as tool-using
agents created an immediate demand for a standardized
context-injection contract, which MCP filled. Second, multi-agent
deployment patterns made it apparent that a tool-binding protocol
alone was insufficient. Those patterns include supervisor-and-worker,
peer-to-peer specialist teams, and long-running task delegation
across organizational boundaries. Peer-messaging protocols (A2A,
ACP, ANP, AGNTCY, Coral) emerged within months of MCP's
stabilization. Third, the
rise of agent-initiated transactions on behalf of human principals
exposed an unmet need for cryptographically verifiable
authorization, leading the major card networks and a parallel
on-chain effort (ERC-8004) to specify what we treat in this survey
as the commerce-settlement layer of inter-agent coordination.

The result is a landscape with substantial overlap but no clear
ranking: a practitioner choosing among A2A, ANP, and Coral for a
multi-agent deployment finds extensive specifications but no
neutral basis for comparison. The four prior surveys provide
taxonomies and qualitative observations; this survey provides a
normalized scoring framework, anchored against the eight
protocols of the tool-binding and peer-messaging layers and
applied to thirteen protocols in total across the three layers,
alongside an industry
adoption synthesis and a security threat-coverage matrix.

\subsection{Contributions}
\label{sec:intro:contributions}

This survey makes seven distinct contributions.
Contributions C1 to C4 are survey contributions, each of
which is absent from the four prior surveys of the
landscape~\cite{yang2025survey, ehtesham2025interop,
kong2025communication, commcentric2025survey}.
Contributions C5, C6, and C7 are protocol-design
proposals derived from the open-problems agenda the
survey identifies; they are narrowly-scoped extensions
to the contemporary protocol stack, not free-standing
protocols, and they are presented in
\S\ref{sec:proposals} after the survey is complete.

\paragraph{Contribution C1: A uniform 10-axis rubric.}
We propose an evaluation rubric of ten axes. The axes are Discovery
and Registry
Model, Identity and Authentication, Authorization and Delegation,
Transport and Wire Format, Task and Conversation Model, Async and
Long-Running Operations, Multi-Modal Content Handling, Security
Threat Surface, Semantic Interoperability, and Governance / License /
Observability. We apply the ten axes to the eight tool-binding and
peer-messaging protocols (\S\ref{sec:rubric};
Table~\ref{tab:rubric-matrix}), and compare the four
commerce-settlement protocols on the structural axes that do apply to a
settlement envelope (Table~\ref{tab:commerce-stack}). Every rubric score
in this survey is therefore one of eight protocols, and the count of
thirteen refers to the surveyed set. Seven axes use
anchored ordinal scoring (0 to 3) with explicitly named anchor
descriptors per level; one axis (Transport) is a structured
tuple to avoid the false precision of ordinal scoring over a
configuration space; one axis (Security) is a count over a fixed
twelve-item threat catalog; and one axis (Governance / License /
Observability) is a composite triple of categorical steward,
license identifier, and 0 to 3 observability score. Prior surveys propose taxonomies
(Yang~\cite{yang2025survey}, Kong~\cite{kong2025communication}) or
adoption roadmaps (Ehtesham~\cite{ehtesham2025interop}) but no
normalized scoring framework that lets a reader rank a new protocol
against existing ones.

\paragraph{Contribution C2: First-class treatment of the
agentic-commerce stack.}
We treat AP2, Visa TAP, Mastercard Agent Pay, ERC-8004, and UCP as a
first-class third architectural layer of inter-agent coordination,
alongside tool-binding (MCP and its descendants) and peer-messaging
(A2A, ANP, ACP, AGNTCY, Coral, LOKA, ACNBP). Section~\ref{sec:industry-protocols}
analyzes the commerce stack at the spec level. That analysis covers
its intent-mandate-receipt envelope, its principal-delegation chain
semantics, its on-chain versus off-chain settlement models, and
the cross-network interoperability work underway in the Agentic
AI Interoperability Foundation~\cite{lf_aaif}. It also covers UCP's
alternative to that harmonisation: fix the transaction layer above
the settlement protocols and treat the mandate as a negotiated
extension. No prior survey
covers the commerce-settlement layer at this depth; Yang and
Kong predate its emergence; Ehtesham's four-protocol scope omits
it; the communication-centric survey~\cite{commcentric2025survey}
treats payments as out of scope.

\paragraph{Contribution C3: Industry adoption mapped by vertical.}
Section~\ref{sec:adoption} and Table~\ref{tab:industry-adoption}
synthesize industry adoption across four verticals: enterprise
SaaS, agentic commerce and payments, healthcare, and finance. The
synthesis names production deployments rather than treating protocols
in the abstract. We track which protocols are shipping in identifiable
products, which carry pilot-stage announcements, and which remain
research artifacts. Prior surveys do not provide this mapping; the
absence is particularly notable in Yang and Kong, where adoption
is treated only at the level of which vendor designed which
specification.

\paragraph{Contribution C4: A machine-readable companion artifact.}
We ship a companion repository (\companionrepo) carrying \texttt{protocol-scores.yaml}
(the source of truth behind Table~\ref{tab:rubric-matrix},
Figure~\ref{fig:heatmap}, and Figure~\ref{fig:radar}), the threat
coverage and adoption matrices as YAML, and the scripts that
regenerate the tables and figures built from them. This artifact addresses the
staleness risk that broadly threatens any protocol survey across
a 6 to 9 month review cycle: scores remain re-runnable as protocol
versions drift, and a reader can fork the YAML to score a ninth
protocol without re-implementing the comparison. Prior surveys do
not provide machine-readable comparison artifacts. Its contents
and license terms are listed in the reproducibility checklist of
Appendix~\ref{sec:appendix:repro}.

\paragraph{Contribution C5: The OWA scoping scheme.}
We propose a three-tier \emph{org / workspace / agent} URI naming
scheme (\S\ref{sec:proposals:owa}) that extends the
\texttt{agent://} URI scheme of Rodriguez~\cite{agent_uri_scheme}
with an explicit workspace tier between the organizational trust
root and the capability path. The workspace tier is missing from
every contemporary identifier scheme: \texttt{agent://}, the
Agent Name Service IETF draft~\cite{ans_ietf}, A2A AgentCard URIs,
W3C DIDs, and Ruflo's cross-machine federation~\cite{ruflo}.
Yet the empirical reality of enterprise deployments points the
other way (\S\ref{sec:adoption:saas}). Multiple agents with
overlapping capabilities operate under different workspace
governance regimes. OWA names the workspace as a first-class
URI segment, enabling per-workspace policy enforcement, cross-org
trust bootstrap through DNS, and the boundary detection that
Contribution C6 depends on. The proposal addresses P7
(trust-bootstrapping across heterogeneous registries) of the
open-problems agenda (\S\ref{sec:openproblems:trust}).

\paragraph{Contribution C6: The AGRP guardrails envelope.}
We propose a normative envelope extension to A2A
(\S\ref{sec:proposals:agrp}) that wraps every cross-boundary
inter-agent message with classification labels, a
fingerprint-evidenced redaction log, and cryptographic
attestation that an egress filter ran. The envelope is bound to
the sender's OWA identity (Contribution C5), allowing the
receiver to verify the egress-filtering property at the
\emph{message} layer rather than relying on the sending
deployment to enforce it. The proposal is positioned explicitly
as a composition partner with the most-developed deployment-layer
precedent, Pipelock~\cite{pipelock}: the AGRP envelope's
\texttt{mediator\_receipt} field is designed to carry Pipelock
receipts forward, yielding a defence-in-depth deployment of
deployment-layer enforcement plus protocol-layer cross-vendor
verifiability. AGRP addresses three open security problems, as
catalogued in \S\ref{sec:security:open}. They are
O3 (prompt-injection mitigation as spec-level requirement),
O5 (contextual integrity for forwarding), and O6 (side-channel
metadata privacy).

\paragraph{Contribution C7: The ACSP context selection envelope.}
We propose a sibling envelope extension to A2A
(\S\ref{sec:proposals:acsp}) that records the
\emph{declared selection step itself}: which candidate items from
the sender's context pool were considered, which were
kept, which were dropped, under which token budget and
optimisation objective, with a keyed per-candidate
commitment and an attestation chain signed by the
sender's OWA-bound key (Contribution C5). Where AGRP
attests the cleaning of \emph{spans within} forwarded
items, ACSP attests the prior decision of \emph{which
items} crossed the boundary. That is the gap AgentLeak's
output-only audits miss 41.7\% of the
time~\cite{agentleak}. ACSP is positioned as the
protocol-layer answer to the contextual-integrity
question \S\ref{sec:security:privacy} leaves open, and
is anchored on the PACMS~\cite{pacms_cikm_demo} reference
implementation (submodular facility-location selection;
the classical coverage guarantee is not asserted for the
audited token-knapsack implementation without its required
conditions). The three proposals are co-designed: OWA
names the boundary; ACSP selects which items cross it;
AGRP redacts the spans within those items. The triplet
closes the contextual-integrity loop at the wire layer
in a way no contemporary protocol provides.

\subsection{What this survey is not}
\label{sec:intro:notwhat}

The scope of this survey is deliberately bounded, and several
adjacent literatures are explicitly out of frame.

This is \emph{not} a survey of LLM agents or single-agent reasoning
frameworks. AutoGen~\cite{wu2023autogen},
MetaGPT~\cite{hong2024metagpt}, ChatDev~\cite{qian2024chatdev},
CAMEL~\cite{li2023camel}, and their successors
are foundational context for the protocols we evaluate, but they
are not themselves subjects of the rubric and do not appear in
Table~\ref{tab:rubric-matrix}.

This is \emph{not} a survey of multi-agent reinforcement learning.
Where coordination questions overlap (e.g., consensus, role
assignment), we cite the MARL literature for context only.

This is \emph{not} a re-derivation of the security threat
literature. The deep security analysis of inter-agent communication
already exists in Kong et al.~\cite{kong2025communication} and the
parallel red-team
catalogs~\cite{a2asecbench, zhan2024injecagent, a2a_weaknesses}.
Section~\ref{sec:security} synthesizes these as a 12-threat
spec-level coverage matrix per protocol; for primary security
analysis the reader is referred to Kong directly.

This is \emph{predominantly} a survey, not a protocol-design
proposal. The rubric (Contribution C1) scores existing
specifications; it does not specify a new one. The three
exceptions are Contributions C5, C6, and C7, which sit in
a final dedicated section (\S\ref{sec:proposals}) and are
derived directly from the open-problems agenda
(\S\ref{sec:openproblems}). The proposals are narrowly
scoped: a workspace-aware URI extension to the
\texttt{agent://} scheme~\cite{agent_uri_scheme}, a
normative guardrails envelope extension to A2A, and a
sibling context-selection envelope. They are
positioned as candidates the AAIF and W3C working groups
can adopt, extend, or reject. They are reference designs,
not validated specifications; \S\ref{sec:proposals} is
explicit about the implementation-maturity caveat.

Finally, the survey is bounded by a declared cutoff date of
\textbf{2026-Q1}. Specifications and primary literature published
after that date may be referenced in footnotes but are not
exhaustively covered. The companion repository carries the
canonical updatable rubric scores past that horizon.

\subsection{Reading guide}
\label{sec:intro:guide}

Section~\ref{sec:background} locates the 2024 to 2026 protocol wave
in the prior heritage of agent communication languages, from
FIPA-ACL and KQML through SOAP/WSDL, REST, and OpenAPI to the
JSON-RPC 2.0 substrate that all current inter-agent protocols
share. Section~\ref{sec:taxonomy} introduces the three-layer
architectural grid (tool-binding, peer-messaging,
commerce-settlement) and the orthogonal classification axes
(open versus consortium-governed; registry-based versus
peer-discovery; sync versus async-first) that organize the
remainder of the survey. Section~\ref{sec:rubric} defines the
ten-axis rubric in full; it is the novelty anchor of the survey
and is the only section the time-pressed reader should not skim.
Sections~\ref{sec:industry-protocols} to \ref{sec:emerging-protocols}
apply the rubric to industry-led and research-led protocols
respectively. Section~\ref{sec:comparison} presents the master
rubric matrix, the heatmap, and the radar charts.
Section~\ref{sec:security} synthesizes the security threat
literature into a 12-threat coverage matrix.
Section~\ref{sec:semantic} treats semantic interoperability as
a distinct concern from message passing.
Section~\ref{sec:adoption} maps industry adoption by vertical.
Section~\ref{sec:governance} analyzes governance, licensing, and
observability as institutional factors that decide protocol
survival. Section~\ref{sec:openproblems} catalogues the
unresolved problems that should drive the next standardization
cycle. Section~\ref{sec:proposals} presents the three
narrowly-scoped protocol-design proposals (Contributions
C5, C6, and C7) that derive from the open-problems
agenda, with the three differentiation tables that
position each proposal against its prior art.
Section~\ref{sec:conclusion} concludes.

The reader who wants only the bottom line of the comparison
should turn directly to Table~\ref{tab:rubric-matrix} on
page~\pageref{tab:rubric-matrix} and Figures~\ref{fig:heatmap}
and~\ref{fig:radar}; everything else in the survey is
explanation, evidence, and context for those three artifacts.
